\section{Introduction}

Reciprocal Recommender System (RRS)~\cite{pizzato2010recon} has been widely deployed in online platforms, such as online dating~\cite{pizzato2010recon,xia2015reciprocal,neve2019latent} and recruitment~\cite{le2019towards,jiang2020learning,yang2022modeling}.
Different from traditional recommendation algorithms, it emphasizes a two-way (or bilateral) selection process when establishing the matching between the two parties involved in recommender systems. 
Generally, existing RRS works can be roughly divided into content-based methods and collaborative filtering-based methods. 
The former mainly model user interests from their
personalized profiles~\cite{pizzato2010recon,luo2020rrcn,yildirim2021bideepfm}, while the latter usually learn static user representations to capture long-term user preference based on collaborative filtering~\cite{xia2015reciprocal,neve2019latent,luo2020motif}. 
However, the interests of users tend to dynamically shift along with alterations in the surrounding environment over time, which is disregarded by the above two categories of models.  
For instance, in the context of online recruitment, as individuals gain more work experience and qualifications, they are inclined to increase the criteria in job hunting. 
Therefore, it is essential to effectively model the dynamic preference changes in such a two-way preference selection process, thus making more accurate recommendations that can satisfy the needs of both parties.

To model the dynamic interests of users, sequential recommendation methods mainly capture user preference by modeling the   chronological user behavior sequences~\cite{kang2018self,sun2019bert4rec} in a unilateral setting. 
As a typical approach, they utilize the output in the last time step of the sequential recommendation model for matching prediction. 
Different from conventional recommendation algorithms, which are mainly developed based on the single-perspective active selection from a user to items, 
users play  a dual role in reciprocal recommendation, serving not only as the party to actively select but also as the other party to be passively selected.  
Thus, the RRS necessitates the dynamic modeling of both parties, especially in the context of two-way matching, which is likely to result in a more difficult learning problem. 



As the technical approach, we formulate reciprocal recommendation as a distinctive sequence matching task, where the recommendation prediction is performed based on the matching between the behavior sequences of the involved two parties. 
Specifically, we establish the behavior sequences in dual perspectives, 
and then conduct the two-sided matching  by modeling fine-grained sequential semantic interactions.
While, it takes a high time cost to model fine-grained matching, and we further consider employing self-distillation techniques (commonly used in model compression~\cite{zhang2019your,xu2020improving,liu2020fastbert}) to improve model efficiency.  


Inspired by the above motivations, in this paper, we propose \emph{ReSeq}, a \underline{Re}ciprocal \underline{Seq}uential recommendation method.
There are three major technical contributions in our approach.
Firstly, we model bilateral behavior sequences by a transformer network\cite{vaswani2017attention} based on user's  dual-perspective embeddings (\ie \emph{active} and \emph{passive}).
 Secondly, we design a multi-scale matching prediction method to capture fine-grained sequential interactions at both macro and micro levels. 
Finally,  we largely improve the efficiency of the proposed model via micro-to-macro self-distillation,  while ensuring accuracy and effectiveness.

To the best of our knowledge, this is the first time that dynamic behavior sequences of two parties  are considered in reciprocal recommendation, which can achieve both accurate and efficient recommendations. Extensive experiments on five datasets from two reciprocal recommendation scenarios demonstrate the effectiveness of our proposed method.















